\section{Preliminary results on half-bubbles}\label{sec:prelim}
        
        
        	Let $\Sigma^n\subset \R^{n+1}$ be a half-bubble in the sense above. Let us denote by $\Pi_1$ the closed half-space bounded by $\Pi$ that contains the interior of $\Sigma$ and by $\Pi_2$ the other closed half-space bounded by $\Pi$. The terminology \textsl{half-bubble} can be justified as follows: if we reflect $\Sigma$ across $\Pi$ in $\R^{n+1}$ we get a minimal hypersurface without boundary, which a priori is only $C^1$, but a posteriori is $C^{1,\alpha}$ by means of a standard application of De Giorgi-Nash estimates, hence smooth by Schauder theory. Such a minimal hypersurface shall be denoted by $\check{\Sigma}$ and be referred to as the \textsl{double} of $\Sigma$.
        	
        	
        	In relation to the geometric applications we are about to present, we need to compare the Morse index of $\Sigma$ (as a free boundary minimal surface, thus with suitable boundary conditions) and the Morse index of $\check{\Sigma}$. In fact, our results will follow as a specification of a more general discussion.
        
     
        	
        	\subsection{Schr\"odinger-type operators on involutive manifolds}\label{subs:general}
        	
        	Let $(\check{\Sigma}^{n},g)$ be a complete Riemannian manifold, without boundary (in Section \ref{subs:morsehalf} we will then specify our discussion to the ambient manifold $\check{\Sigma}$ presented in the previous paragraph, namely a bubble with its induced metric). Suppose there exists a Riemannian involution $\tau:(\check{\Sigma},g) \to (\check{\Sigma},g)$, namely a smooth isometry satisfying the two conditions:
        	\[
        	\tau\circ\tau = \operatorname{id}, \ \ \tau\neq \operatorname{id}.
        	\] 
        	
        Given a linear functional space $X$ consisting of functions defined on $(\check{\Sigma},g)$ let us then introduce the subspaces of even and odd functions with respect to the action of $\tau$:
        	\[
        	X_{\jepPubScript{E}}=\left\{\varphi\in X: \varphi\circ\tau=\varphi \right\}, \ X_{\jepPubScript{O}}=\left\{\varphi\in X: \varphi\circ\tau=-\varphi \right\}.
        	\]	
        	In particular, notice that 
        	\[
        	C^{\infty}=C^{\infty}_{{\jepPubScript{E}}}\oplus C^{\infty}_{{\jepPubScript{O}}}.
        	\]
        	This direct sum is orthogonal in $L^2$ when restricted to the subspace of smooth, compactly supported functions. In fact, we have
        	\[
        	L^2=L^2_{{\jepPubScript{E}}}\oplus^{\perp} L^2_{{\jepPubScript{O}}}.
        	\]
        Given a function $V\in C^{\infty}_{{\jepPubScript{E}}}$ we wish to study a Schr\"odinger-type operator of the form
        	\[
        	T_V\varphi=\Delta_g\varphi+V\varphi, 
        	\]
        	together with the associated quadratic form
        	\[
        	Q_V(\varphi,\varphi)=-\int_{\check{\Sigma}}\varphi T_V \varphi\,d\h^n=\int_{\check{\Sigma}}(|\nabla \varphi|^2-V\varphi^2)\,d\h^n,
        	\]
        	which a priori is defined on the set of smooth, compactly supported functions.
        	
        	
        	\begin{definition}\label{def:index} In the setting above, we define
        		\[
        		\operatorname{Ind} (Q_V), \ \ \operatorname{Ind}_{{\jepPubScript{E}}} (Q_V), \ \ \operatorname{Ind}_{{\jepPubScript{O}}} (Q_V)
        		\]
        		as the largest dimension of a linear subspace of 
        		\[
        		C^{\infty}_c, \ \ (C^{\infty}_c)_{{\jepPubScript{E}}}, \ \ (C^{\infty}_c)_{{\jepPubScript{O}}}
        		\]
        		respectively, where the quadratic form $Q_V$ is negative definite.
        	\end{definition}	
        	
        	
        	Thereby, the following inequality is straightforward:
        	
        	\begin{lem}\label{lem:ineq}
        		In the setting above, 
        		\[
        		\operatorname{Ind} (Q_V) \geqslant \operatorname{Ind}_{{\jepPubScript{E}}} (Q_V)+\operatorname{Ind}_{{\jepPubScript{O}}} (Q_V)
        		\]
        		and
        		\[
        		\operatorname{Ind} (Q_V)=0 \iff \operatorname{Ind}_{{\jepPubScript{E}}} (Q_V)=0 \ \text{and} \ \operatorname{Ind}_{{\jepPubScript{O}}} (Q_V)=0.
        		\] 
        	\end{lem}
        	
        	\begin{proof} Since $V$ is even and $\tau$ is an isometry, it is immediate to check that the operator $T_V$ preserves the decomposition of $C^{\infty}$ into even and odd functions. Thus, if $\varphi_{{\jepPubScript{E}}}\in (C^{\infty}_c)_{{\jepPubScript{E}}}$ and $\varphi_{{\jepPubScript{O}}}\in (C^{\infty}_c)_{{\jepPubScript{O}}}$ then
        		\begin{equation}\label{eq:mixedprod}
        		Q_{V}(\varphi_{{\jepPubScript{E}}},\varphi_{{\jepPubScript{O}}})=0,
        		\end{equation}
        		and hence, by bilinearity, given any $\varphi\in C^{\infty}_c$ and writing $\varphi=\varphi_{{\jepPubScript{E}}}+\varphi_{{\jepPubScript{O}}}$ one has
        		\[
        		Q_{V}(\varphi,\varphi)=Q_{V}(\varphi_{{\jepPubScript{E}}},\varphi_{{\jepPubScript{E}}})+Q_{V}(\varphi_{{\jepPubScript{O}}},\varphi_{{\jepPubScript{O}}}),
        		\]
        		which easily implies the claims.
        	\end{proof}		
        	
        	
        	
        	In fact, we claim that the inequality above is actually an equality. 
        	With that goal in mind, we restrict to the case of \textsl{finite} Morse index and recall a simple but useful result:
        	
        	\begin{prop}\label{pro:l2spec}(cf. \cite[Prop.~2]{FC85})
        		In the setting above, the following two statements are equivalent:
        		\begin{enumerate}
        			\item{the index of the quadratic form $Q_V$ is finite;}
        			\item{there exists a finite dimensional subspace $W$ of $L^2$ having an orthonormal basis $\varphi^1,\ldots,\varphi^k$ consisting of eigenfunctions of $T_V$ with eigenvalues $\lambda_1,\ldots,\lambda_k$ respectively. Each $\lambda_i$ is negative and $Q_V(\varphi,\varphi)\geqslant 0$ whenever $\varphi\in C^{\infty}_c\cap W^{\perp}$. In this case $k$ equals the Morse index of $Q_V$.}	
        		\end{enumerate}		
        	\end{prop}
        	
        	
        	\begin{rmk}\label{rmk:Vbounded} If $V$ is assumed to be bounded, then the basis provided in Proposition \ref{pro:l2spec} actually consists of elements of finite Dirichlet energy, i.e. functions belonging to $W^{1,2}$, the closure of $C^{\infty}_c$ with respect to the Sobolev norm determined by
        		\begin{equation}\label{eq:SobNorm}
        		\|\varphi\|^2_{W^{1,2}}:= \|\varphi\|^2_{L^2}+\|\nabla \varphi\|^2_{L^2}.
        		\end{equation}	
        	\end{rmk}	
        	
        	
        	
        	Based on this fact and in view of the geometric applications we wish to present, we \emph{assume} from now onwards that the function	$V\in C^{\infty}_{{\jepPubScript{E}}}$ is uniformly bounded. 
        
        	
        	It is convenient to introduce the relevant notion of (point) spectrum in this setting.
        	
        	\begin{definition}\label{def:spec} In the setting above, we define
        		\[
        		\operatorname{spec} (Q_V):=\Big\{\text{critical values of the map} \ \varphi\in W^{1,2}\smallsetminus \left\{0\right\} \mapsto \frac{Q_V(\varphi,\varphi)}{\|\varphi\|^2_{L^2}}\Big\};
        		\]
        		\[
        		\operatorname{spec}_{{\jepPubScript{E}}} (Q_V):=\Big\{\text{critical values of the map} \ \varphi\in (W^{1,2})_{{\jepPubScript{E}}}\smallsetminus \left\{0\right\} \mapsto \frac{Q_V(\varphi,\varphi)}{\|\varphi\|^2_{L^2}}\Big\};
        		\]
        		\[
        		\operatorname{spec}_{{\jepPubScript{O}}} (Q_V):=\Big\{\text{critical values of the map} \ \varphi\in (W^{1,2})_{{\jepPubScript{O}}}\smallsetminus \left\{0\right\} \mapsto \frac{Q_V(\varphi,\varphi)}{\|\varphi\|^2_{L^2}}\Big\}.
        		\]
        	\end{definition}
        To avoid ambiguities: $\lambda\in \operatorname{spec}(Q_V)$ if and only if we can find an associated critical point $\varphi \in W^{1,2}\smallsetminus\left\{0\right\}$, which satisfies
        	\begin{equation}\label{eq:weak}
        	\int_{\check{\Sigma}}(\nabla \varphi \cdot \nabla\zeta -V\varphi\zeta)\,d\h^n = \lambda \int_{\check{\Sigma}}\varphi\zeta\,d\h^n, \ \ \ \forall \zeta\in C^{\infty}_c,
        	\end{equation}
        	whence it is standard to get that $\varphi$ is actually smooth and solves, in a classical sense, the eigenvalue equation
        	\[
        	\Delta_g \varphi+V\varphi+\lambda\varphi=0
        	\]
        	which also implies that $\Delta_g\varphi \in L^2$. 
        	
        	Similar arguments can be applied to the other two cases as well, with an important caveat. By definition, we have $\lambda\in \operatorname{spec}_{{\jepPubScript{E}}} (Q_V)$ (respectively $\lambda\in \operatorname{spec}_{{\jepPubScript{O}}} (Q_V)$) if equation \eqref{eq:weak} is satisfied for every \textsl{even} (respectively \textsl{odd}) test function $\zeta\in C^{\infty}_c$. However, we also have by symmetry arguments (cf. equation \eqref{eq:mixedprod}) that the same equation is satisfied for every \textsl{odd} (respectively \textsl{even}) test function, and hence for any $\zeta\in C^{\infty}_c$ (in either of the two cases). Therefore, we shall actually conclude
        	\[
        	\lambda\in \operatorname{spec}_{{\jepPubScript{E}}} (Q_V) \ \iff \ \exists \varphi\in  C^{\infty}_{{\jepPubScript{E}}}\smallsetminus \left\{0\right\} \ \text{such that} \	\Delta_g \varphi+V\varphi+\lambda\varphi=0,
        	\]
        	and
        	\[
        	\lambda\in \operatorname{spec}_{{\jepPubScript{O}}} (Q_V) \ \iff \ \exists \varphi\in  C^{\infty}_{{\jepPubScript{O}}}\smallsetminus \left\{0\right\} \ \text{such that} \	\Delta_g \varphi+V\varphi+\lambda\varphi=0.
        	\]
        	
        	
        	
        	
        	\begin{lem}\label{lem:specadd}
        		In the setting above,
        		\[
        		\operatorname{spec} (Q_V)=\operatorname{spec}_{{\jepPubScript{E}}} (Q_V) \cup 	\operatorname{spec}_{{\jepPubScript{O}}} (Q_V).
        		\]
        		Moreover, under the assumption that the operator $T_V$ has finite Morse index on $(\check{\Sigma},g)$, we have
        		\[
        		\operatorname{Ind} (Q_V) = \operatorname{Ind}_{{\jepPubScript{E}}} (Q_V)+\operatorname{Ind}_{{\jepPubScript{O}}} (Q_V).
        		\]
        	\end{lem}
        	
        		
        	
        	\begin{proof}
        		For the first claim, notice that the inclusion $\supseteq$ is obvious so let us discuss the other one. Let then $\lambda \in \operatorname{spec} (Q_V)$ and write $\varphi=\varphi_{{\jepPubScript{E}}}+\varphi_{{\jepPubScript{O}}}$ (the $L^2$ decomposition of an associated eigenfunction into even and odd parts), thus by linearity
        		\[
        		0=\Delta_g \varphi+V\varphi+\lambda\varphi= \underbrace{(\Delta_g \varphi_{{\jepPubScript{E}}}+V\varphi_{{\jepPubScript{E}}}+\lambda\varphi_{{\jepPubScript{E}}})}_{\in L^2_{{\jepPubScript{E}}} } + \underbrace{(\Delta_g \varphi_{{\jepPubScript{O}}}+V\varphi_{{\jepPubScript{O}}}+\lambda\varphi_{{\jepPubScript{O}}})}_{\in L^2_{{\jepPubScript{O}}}}
        		\]
        		hence each of the two summands must vanish and thus (since either $\varphi_{{\jepPubScript{E}}}\neq 0$ or $\varphi_{{\jepPubScript{O}}}\neq 0$) we get $\lambda\in \operatorname{spec}_{{\jepPubScript{E}}} (Q_V)$ or $\lambda\in \operatorname{spec}_{{\jepPubScript{O}}} (Q_V)$ as it was claimed.
        		
        		
        		
        		For the second part, recall that by Proposition \ref{pro:l2spec} there exist (under the assumption that $T_V$ has finite Morse index) finitely many (say $k$) eigenfunctions $\varphi^1,\ldots,\varphi^k \in L^2$ that correspond to the negative eigenvalues $\lambda_1,\ldots,\lambda_k$ (where it is understood that each eigenvalue can be repeated if it comes with multiplicity). 
        		Following the argument we just presented, replace each $\varphi^j$ by means of the couple $\varphi^j_{{\jepPubScript{E}}}, \varphi^j_{{\jepPubScript{O}}}$ and set
        		\[
        		V=\operatorname{span}_{\R}\left\{\varphi^1,\ldots,\varphi^k \right\}, \ \ \widetilde{V}=\operatorname{span}_{\R}\left\{\varphi^1_{{\jepPubScript{E}}}, \varphi^1_{{\jepPubScript{O}}},\ldots, \varphi^k_{{\jepPubScript{E}}}, \varphi^k_{{\jepPubScript{O}}} \right\}.
        		\] 
        		Now, since $\varphi^j=\varphi^j_{{\jepPubScript{E}}}+\varphi^j_{{\jepPubScript{O}}}$ for any $j\in\left\{1,\ldots, k\right\}$ we have
        		\[
        		\text{dim}_{\R}(\widetilde{V})\geqslant k
        		\]
        		and thus there are at least $k$ functions in the collection $\left\{\varphi^1_{{\jepPubScript{E}}}, \varphi^1_{{\jepPubScript{O}}},\ldots, \varphi^k_{{\jepPubScript{E}}}, \varphi^k_{{\jepPubScript{O}}} \right\}$ which are not zero, and eigenfunctions either for $\operatorname{spec}_{{\jepPubScript{E}}} (Q_V)$  or for $\operatorname{spec}_{{\jepPubScript{O}}} (Q_V)$, with negative eigenvalues. Thereby, the inequality one gets, combined together with Lemma \ref{lem:ineq}, enables us to complete the proof.
        	\end{proof}	
        	
        	
        	
